% Generated draft: requires mathematical and source review.
\documentclass[11pt]{article}
\usepackage[margin=25mm]{geometry}
\usepackage{fontspec}
\setmainfont{Latin Modern Roman}
\usepackage{amsmath,amssymb,mathtools,unicode-math}
\setmathfont{Latin Modern Math}
\usepackage{enumitem,needspace,xurl,hyperref}
\setlength{\parindent}{0pt}
\setlength{\parskip}{5pt}
\setlength{\emergencystretch}{3em}
\newcommand{\N}{\mathbb{N}}
\newcommand{\Z}{\mathbb{Z}}
\newcommand{\Q}{\mathbb{Q}}
\newcommand{\R}{\mathbb{R}}
\newcommand{\C}{\mathbb{C}}
\newcommand{\F}{\mathbb{F}}
\newcommand{\A}{\mathbb{A}}
\newcommand{\PP}{\mathbb{P}}
\newcommand{\E}{\mathbb{E}}
\newcommand{\eps}{\varepsilon}
\newcommand{\dd}{\,\mathrm{d}}
\begin{document}
\section*{UnsolvedMath ID: 30004730}
\textbf{Source problem code:} catalogue entry
\subsection{Definitions and mathematical statement}
Let \((X,d)\) be a metric space. A geodesic bicombing on \(X\) is a map
\[
\sigma:X\times X\times[0,1]\longrightarrow X
\]
such that, for every \(x,y\in X\) and \(s,t\in[0,1]\),
\[
\sigma(x,y,0)=x,\qquad \sigma(x,y,1)=y,\qquad
 d\bigl(\sigma(x,y,s),\sigma(x,y,t)\bigr)=|s-t|\,d(x,y).
\]
It is conical if, for every \(x,y,x',y'\in X\) and \(t\in[0,1]\),
\[
 d\bigl(\sigma(x,y,t),\sigma(x',y',t)\bigr)
 \leq (1-t)d(x,x')+t\,d(y,y').
\]
It is consistent if, for every \(x,y\in X\), \(0\leq a\leq b\leq1\), and \(u\in[0,1]\),
\[
\sigma\bigl(\sigma(x,y,a),\sigma(x,y,b),u\bigr)
 =\sigma\bigl(x,y,(1-u)a+ub\bigr).
\]
These definitions use directed segments; reversibility, the additional identity \(\sigma(y,x,t)=\sigma(x,y,1-t)\), is not imposed. The supplied record does not specify whether its terminology includes reversibility, so that convention requires source confirmation.

\textbf{Question as recorded.} Is it true that, for every metric space \((X,d)\), the existence of a conical geodesic bicombing on \(X\) implies the existence of a consistent conical geodesic bicombing on \(X\)? The latter bicombing may differ from the initially given one.

\textbf{Unresolved scope qualification.} The recorded question imposes no completeness hypothesis, whereas the supplied literature assessment describes the unresolved general assertion for complete metric spaces. A metric space is complete if every Cauchy sequence converges in it. The record does not resolve whether completeness was intended as a hypothesis; the question above retains the literal quantification over all metric spaces.

The supplied dated curation labels the problem open, reports that the proper-space case is settled under unspecified weaker convexity assumptions, and reports no known counterexample to the complete-space assertion. Here proper means that every closed bounded subset is compact. These are supplied literature assessments, not independently verified mathematical claims; the unspecified weaker assumptions do not determine a precise auxiliary theorem.
\subsection{Short English statement}
If a metric space admits a conically controlled choice of geodesic between every pair of points, must it admit such a choice that agrees with its own subsegments? The record leaves completeness and the convention concerning reversibility unresolved. Its open-status assessment is supplied dated curation.
\subsection{Sources}
\begin{itemize}
\item Local UnsolvedMath registry, supplied complete source record: permanent ID \(30004730\), title \emph{Consistent Conical Bicombings in Metric Spaces}, exact source problem code \texttt{catalogue entry}, in the Oberwolfach Reports Open Problems collection. The record gives proposed year \(2021\) and names no proposer. Its supplied status is \texttt{open}; its statement status is \texttt{corrected\_verified}, with a supplied claim of manual semantic, typographic, and LaTeX review. These labels and that review claim have not been independently verified here.
\item Supplied report citation: \emph{Differentialgeometrie im Grossen (2022)}, Oberwolfach Reports, pp.\ 1685--1734, DOI \texttt{10.4171/owr/2021/32}. \url{https://doi.org/10.4171/owr/2021/32}. The record describes Oberwolfach Reports as documenting mathematical workshops and their open-problem discussions. The citation's parenthetical year \(2022\) differs from the DOI's year component and the recorded proposed year \(2021\); this discrepancy is unresolved.
\item Supplied literature source: ``Basso, conical bicombings.'' \url{https://ems.press/journals/lem/articles/14297776}. No fuller bibliographic details are supplied.
\item Supplied literature source: ``Metric geometry paper on consistent bicombings.'' \url{https://doi.org/10.1007/s00208-024-02905-w}. No authors or precise title are supplied.
\item Literature assessment checked at \texttt{2026-08-22}, exactly as recorded. This is supplied dated curation derived from the supplied OWR curation, not an independently verified assessment. It labels the problem open, reports a result for proper spaces under weaker convexity assumptions, and describes the complete-space assertion as open with no known counterexample. The record itself requests independent verification of time-sensitive status and cited claims.
\end{itemize}
\end{document}
